\documentclass[12pt]{article}
\usepackage[T1]{fontenc}
\usepackage[utf8]{inputenc}
\usepackage{lmodern}
\usepackage{textcomp}
\usepackage{enumitem}
\usepackage{geometry}
\usepackage{graphicx}
\usepackage{amsmath, amssymb}
\geometry{margin=1in}
\begin{document}
\begin{center}
Chemistry - Graduate Level Assessment 4 \
Total Marks: 40 \
Answer all questions.
\end{center}

\section*{Multiple Choice (10 marks)}
\begin{enumerate}[label=\arabic*.]
\item Evaluate the series $ S=\frac{1}{2}+\frac{1}{4}+\frac{1}{8}+\frac{1}{16}+\cdots $\$
\begin{enumerate}[label=(\alph*)]
    \item 0.5
    \item 2.5
    \item 1.25
    \item 1.5
    \item 4
\end{enumerate}
\item A chemist wants to produce chlorine from molten NaCl. How many grams could be produced if he uses a current of one amp for 5.00 minutes?
\begin{enumerate}[label=(\alph*)]
    \item 0.031 g
    \item 5.55 g
    \item 0.021 g
    \item 0.511 g
    \item 0.111 g
\end{enumerate}
\item 2.000 picogram(pg) of ^33 P decays by ^0\_-1\ensuremath{\beta} emission to 0.250 pg in 75.9 days. Find the half-life of ^33 P.
\begin{enumerate}[label=(\alph*)]
    \item 37.95 days
    \item 50.6 days
    \item 8 days
    \item 5.6 days
    \item 60.4 days
\end{enumerate}
\item The photoelectron spectrum of carbon has three equally sized peaks. What peak is at the lowest energy?
\begin{enumerate}[label=(\alph*)]
    \item The 2 p peak has the lowest energy.
    \item The 3 d peak has the lowest energy.
    \item The 2 d peak has the lowest energy.
    \item The 1 d peak has the lowest energy.
    \item The 1 p peak has the lowest energy.
\end{enumerate}
\item Calculate the voltage of the cell Fe; Fe^+2 \ensuremath{\vert}\ensuremath{\vert} H^+ ; H\_2 if the iron half cell is at standard conditions but the H^+ ion concentrations is .001 M.
\begin{enumerate}[label=(\alph*)]
    \item 0.50 volts
    \item 0.76 volts
    \item 1.2 volts
    \item 0.33 volts
    \item .85 volt
\end{enumerate}
\end{enumerate}

\section*{True / False (10 marks)}
\textit{Answer True or False}
\begin{enumerate}[label=\arabic*.]
\setcounter{enumi}{5}
\item True or False: Thionyl chloride (SOCl 2) has a tetrahedral molecular geometry due to its four bonded atoms around the sulfur atom.
\item True or False: The Q-factor for an X-band EPR cavity with a resonator bandwidth of 1.58 MHz is Q = 5012.
\item True or False: The value of (0.$0081)^(3$/4) is equal to (1000) / (81).
\item True or False: To prepare a 50 ml solution of pH 1.25 HCl, one must add 49.76 ml of H\_2 O to 0.24 ml of concentrated HCl.
\item True or False: The molarity of the sodium hydroxide solution is 0.0641 M after neutralization with hydrochloric acid.
\end{enumerate}

\section*{Long Form Response (20 marks)}
\textit{Answer each question with a clear and complete explanation.}
\begin{enumerate}[label=\arabic*.]
\setcounter{enumi}{10}
\item Calculate the pH of a solution formed by mixing 50.0 mL of 0.0025 M HBr with 50.0 mL of 0.0023 M KOH. Discuss the steps involved in your calculation and the underlying principles.
\item A particle of mass $9.1 \times 10^{-28} \mathrm{~g}$ in a one-dimensional quantum box transitions from the $n=5$ to $n=2$ level, emitting a photon of frequency $6.0 \times 10^{14} \mathrm{~s}^{-1}$. Calculate the length of the box and discuss the significance of this length in the context of quantum confinement.
\end{enumerate}

\vfill
\noindent\textit{End of Paper}
\end{document}